\documentclass[11pt]{article}
\usepackage{amsmath,amssymb,amsthm}
\usepackage[margin=1.1in]{geometry}
\usepackage{booktabs}
\newtheorem{proposition}{Proposition}
\theoremstyle{remark}
\newtheorem{remark}{Remark}

\title{Multiclass SVM:\\ Crammer--Singer versus One-vs-Rest}
\author{Yuval Lavie}
\date{\today}

\begin{document}
\maketitle

\begin{abstract}
Two lifts of the soft-margin SVM to $K$ classes: the one-vs-rest
reduction ($K$ independent binary problems) and the Crammer--Singer
joint objective, whose margin separates the true class's score from
the \emph{best wrong} one. The joint subgradient is derived and
verified against automatic differentiation to machine precision, both
methods are trained on identical data, and the comparison is reported
honestly. Companion program: \texttt{multiclass\_svm.py}.
\end{abstract}

\section{The two objectives}

\paragraph{One-vs-rest.} For each class $k$, train a binary
soft-margin SVM with labels $\pm1$ (class $k$ vs.\ the rest); predict
$\arg\max_k \langle w_k, x\rangle$. The $K$ problems never
coordinate, and their scores are not calibrated against each other.

\paragraph{Crammer--Singer.} One regularized problem over
$W \in \mathbb{R}^{K \times d}$:
\begin{equation}
  \min_W \;\; \frac{\lambda}{2}\|W\|_F^2
  + \frac1n \sum_{i=1}^{n}
  \max\Bigl(0,\; 1 + \max_{k \ne y_i} \langle W_k, x_i\rangle
  - \langle W_{y_i}, x_i\rangle\Bigr),
  \label{eq:cs}
\end{equation}
the batch, regularized form of the multiclass hinge of
multiclass-pa.pdf: every row is optimized \emph{against} the others.

\begin{proposition}[Subgradient]
Let $r_i = \arg\max_{k \ne y_i}\langle W_k, x_i\rangle$ and let
$A = \{i : \text{the hinge in \eqref{eq:cs} is positive}\}$. A
subgradient of \eqref{eq:cs} is
\begin{equation}
  \boxed{\;G = \lambda W + \frac1n \sum_{i \in A}
  \bigl(e_{r_i} - e_{y_i}\bigr) x_i^{\!\top}\;}
  \label{eq:grad}
\end{equation}
--- on each violating example, demote the best wrong row and promote
the true one (the perceptron/PA update direction, regularized).
Verified against autograd at a random $W$: maximum entrywise gap
below $10^{-12}$, loss values equal to $10^{-12}$.
\end{proposition}

\section{Numerical comparison}

Three overlapping Gaussian classes ($1100$ train / $400$ test,
seed 7), $30$ subgradient epochs each:

\begin{center}
\begin{tabular}{lll}
\toprule
Claim & Reference & Measured\\
\midrule
Hand subgradient vs.\ autograd & \eqref{eq:grad} & $< 10^{-12}$\\
Objective decreases & \eqref{eq:cs} & monotone over epochs\\
Crammer--Singer accuracy & --- & $0.9775$\\
One-vs-rest accuracy & --- & $0.9825$\\
\bottomrule
\end{tabular}
\end{center}

\begin{remark}[An honest verdict]
On well-separated low-dimensional blobs the reduction matches the
joint objective --- here OvR even edges it within noise. The joint
formulation's advantages are structural, not universal: one
$\lambda$, mutually calibrated scores, a single margin geometry, and
a direct KKT/dual theory (support patterns across classes) --- and
they matter most with many classes and shared structure. Choosing
the fancier objective is an experiment, not a reflex; this repo's
rule since the GPT ablations.
\end{remark}

\end{document}
